
Up to this point, simplex-chains have been treated like traditional blockchains, where each block has only one parent.
Since the vast majority of a simplex-chain's security is provided by other simplex-chains (and only a small proportion comes from that chain's foundational consensus method), are attacks like an empty-block Denial of Service\footnote{
    For an example of this attack, see \citeLJCoiledcoinLink.
} (DoS) possible?
If a simplex-chain were to use PoW, then it might be (relatively) trivial for an attacker to perform such an attack.
This is because --- in traditional blockchains --- controlling more than 50\% of the blocks produced provides \emph{exclusive} control over \emph{which candidate child blocks win} (i.e., are accepted into the canonical chain).\footnote{
    Exclusive control of this nature also allows for protocol changes via soft-forks.
    Additionally, because these soft-forks can be undetectable (when done well) and don't necessarily affect the income of a miner substantially, bribery (of miners) is theoretically inexpensive.
    WRT UT, this problem is resolved in \autoref{sec:preventing-dos-attacks}.
}
Is there a way that we can mitigate this risk?
If blocks were permitted \emph{more} than a single parent, can this \emph{exclusivity} be denied?
